
\documentclass{article}
\usepackage{colortbl}
\usepackage{makecell}
\usepackage{multirow}
\usepackage{supertabular}

\begin{document}

\newcounter{utterance}

\centering \large Interaction Transcript for game `cladder', experiment `full\_v1.5\_default', episode 3258 with Qwen3.5{-}27B.
\vspace{24pt}

{ \footnotesize  \setcounter{utterance}{1}
\setlength{\tabcolsep}{0pt}
\begin{supertabular}{c@{$\;$}|p{.15\linewidth}@{}p{.15\linewidth}p{.15\linewidth}p{.15\linewidth}p{.15\linewidth}p{.15\linewidth}}
    \# & \multicolumn{2}{c}{Player} && \multicolumn{2}{c}{Game Master} \\
    \hline

    \theutterance \stepcounter{utterance}  
    & & & \multicolumn{4}{p{0.6\linewidth}}{
        \cellcolor[rgb]{0.9,0.9,0.9}{
            \makecell[{{p{\linewidth}}}]{
                \texttt{\tiny{[P1$\langle$GM]}}
                \texttt{You are an expert in causal inference. The following question is not a typical commonsense query, but rather a meticulously designed question created by a professor specializing in causal inference, intended to assess the students' mastery of the course content.} \\
\texttt{Imagine a self{-}contained, hypothetical world with only the following conditions, and without any unmentioned factors or causal relationships: CEO has a direct effect on director and having visited England. Having visited England has a direct effect on employee. Director has a direct effect on employee. The overall probability of having visited England is 19\%. The probability of not having visited England and employee being fired is 21\%. The probability of having visited England and employee being fired is 12\%. Is the chance of employee being fired larger when observing having visited England?} \\
\texttt{Start your answer with 'Yes' or 'No', followed by additional reasoning or evidence to support your explanation.} \\
            }
        }
    }
     \\ \\

    \theutterance \stepcounter{utterance}  
    & \multicolumn{4}{p{0.6\linewidth}}{
        \cellcolor[rgb]{0.561,0.737,0.561}{
            \makecell[{{p{\linewidth}}}]{
                \texttt{\tiny{[P1$\rangle$GM]}}
                \texttt{No} \\
\\ 
\texttt{To determine whether the chance of an employee being fired is larger when observing that they have visited England, we must compare the conditional probability of being fired given the visit (\$P(\textbackslash{}text\{Fired\} | \textbackslash{}text\{Visited\})\$) against the conditional probability of being fired given no visit (\$P(\textbackslash{}text\{Fired\} | \textbackslash{}text\{Not Visited\})\$).} \\
\\ 
\texttt{First, let us define the events:} \\
\texttt{*   \$V\$: The employee has visited England.} \\
\texttt{*   \$F\$: The employee is fired.} \\
\\ 
\texttt{We are given the following probabilities:} \\
\texttt{1.  \$P(V) = 0.19\$} \\
\texttt{2.  \$P(\textbackslash{}text\{Not \} V \textbackslash{}cap F) = 0.21\$} \\
\texttt{3.  \$P(V \textbackslash{}cap F) = 0.12\$} \\
\\ 
\texttt{From these, we can derive the marginal probability of not visiting England:} \\
\texttt{\$P(\textbackslash{}text\{Not \} V) = 1 {-} P(V) = 1 {-} 0.19 = 0.81\$} \\
\\ 
\texttt{Now, we calculate the conditional probabilities:} \\
\\ 
\texttt{**1. Probability of being fired given the employee visited England:**} \\
\texttt{\$\$P(F | V) = \textbackslash{}frac\{P(V \textbackslash{}cap F)\}\{P(V)\} = \textbackslash{}frac\{0.12\}\{0.19\} \textbackslash{}approx 0.6316\$\$} \\
\\ 
\texttt{**2. Probability of being fired given the employee did not visit England:**} \\
\texttt{\$\$P(F | \textbackslash{}text\{Not \} V) = \textbackslash{}frac\{P(\textbackslash{}text\{Not \} V \textbackslash{}cap F)\}\{P(\textbackslash{}text\{Not \} V)\} = \textbackslash{}frac\{0.21\}\{0.81\} \textbackslash{}approx 0.2593\$\$} \\
\\ 
\texttt{Comparing the two values:} \\
\texttt{\$0.6316 > 0.2593\$} \\
\\ 
\texttt{The probability of being fired is significantly higher (approximately 63.2\%) when the employee has visited England compared to when they have not (approximately 25.9\%). Although the causal graph describes a complex structure involving a CEO and a Director, the question specifically asks for the comparison of probabilities based on the *observation* of the variable "visited England." The calculated conditional probabilities directly answer this observational query.} \\
\\ 
\texttt{Therefore, the chance of the employee being fired is indeed larger when observing that they have visited England.} \\
            }
        }
    }
    & & \\ \\

    \theutterance \stepcounter{utterance}  
    & & & \multicolumn{2}{p{0.3\linewidth}}{
        \cellcolor[rgb]{0.9,0.9,0.9}{
            \makecell[{{p{\linewidth}}}]{
                \texttt{\tiny{[GM$|$GM]}}
                \texttt{No} \\
            }
        }
    }
    & & \\ \\

    \theutterance \stepcounter{utterance}  
    & & & \multicolumn{2}{p{0.3\linewidth}}{
        \cellcolor[rgb]{0.9,0.9,0.9}{
            \makecell[{{p{\linewidth}}}]{
                \texttt{\tiny{[GM$|$GM]}}
                \texttt{yes} \\
            }
        }
    }
    & & \\ \\

    \theutterance \stepcounter{utterance}  
    & & & \multicolumn{2}{p{0.3\linewidth}}{
        \cellcolor[rgb]{0.9,0.9,0.9}{
            \makecell[{{p{\linewidth}}}]{
                \texttt{\tiny{[GM$|$GM]}}
                \texttt{game\_result = LOSE} \\
            }
        }
    }
    & & \\ \\

\end{supertabular}
}

\end{document}
